% Faithful transcription — original Latin. Transcribed from the original first printing in the
% Novi Commentarii Academiae Scientiarum Imperialis Petropolitanae, tomus VI (St Petersburg,
% 1761), printed pages 37–57: Euler's "De integratione aequationis differentialis
% m dx/√(1−x⁴) = n dy/√(1−y⁴)" (Eneström index E251). Scan: Euler Archive full-text PDF
% (scholarlycommons.pacific.edu/euler-works/251, article=1250).
% \origpage uses the PRINTED page numbers. Per HOUSESTYLE R12, the 1761 orthography is kept
% faithfully: capital V for U (EVLERO, AEQVATIONIS), no u/v distinction (vt, vti, inuoluere),
% author notation (zz for z², xx for x²) — only long-ſ and ligatures are normalized.
\documentclass{article}
\usepackage{readmasters}
\begin{document}

\origpage{37}
\section*{De integratione aequationis differentialis}

\[ \frac{m\,dx}{\sqrt{1-x^{4}}} = \frac{n\,dy}{\sqrt{1-y^{4}}} \]

\emph{Auctore L. EVLERO.}

\textbf{§.~1.} Cum primum occasione inuentionum Ill.~Comitis Fagnani hanc aequationem essem
contemplatus, eiusmodi quidem relationem algebraicam inter variabiles $x$ et $y$ elicui, quae huic
aequationi satisfaceret; sed ea relatio non pro aequatione integrali completa haberi poterat,
propterea quod non complecteretur quantitatem constantem arbitrariam, cuiusmodi semper in calculum
per integrationem introduci solet. Hinc enim, vti satis notum est, integralia incompleta et
particularia distingui solent, quorum illa totam vim aequationum differentialium exhauriunt, haec
vero tantum ita satisfaciunt, vt aliae insuper expressiones aeque satisfacere queant. Criterium
autem aequationis integralis completae in hoc consistit, quod ea quantitatem constantem inuoluere
debeat, quae in aequatione differentiali non apparet.

\origpage{38}
\textbf{§.~2.} Quae quo clarius perspiciantur, sufficiet, aequationem differentialem
simplicissimam $dx = dy$ considerasse, cui vtique satisfacit haec integralis $x = y$, in rem tamen
haec integralis minus late patet, quam differentialis $dx = dy$, cum huic aeque satisfaciat haec
integralis $x = y + a$ multo latius patens, sumendo pro $a$ quantitatem constantem quamcunque,
atque haec demum integralis totam vim aequationis differentialis $dx = dy$ exhaurire censetur, ex
quo etiam aequatio integralis completa appellatur; propterea quod in ea inest quantitas constans
$a$, quae in aequatione differentiali non occurrit. Quodsi vero loco istius constantis indefinitae
$a$ valores determinati substituantur, ex integrali completo obtinentur integralia particularia,
quae ob hanc ipsam rationem minus late patent, quam aequatio differentialis proposita.

\textbf{§.~3.} Saepe numero autem aequationis differentialis integrale particulare algebraicum
exhiberi potest, cum tamen integrale completum sit transcendens; hoc scilicet euenit, si pars
transcendens per constantem illam arbitrariam fuerit multiplicata, quae propterea, constante illa
nihilo aequali posita, ex calculo euanescit, et integrale algebraicum particulare relinquit. Ita
huic aequationi $dy = dx + (y-x)\,dx$ manifestum est, satisfacere valorem $y = x$; quo tamen tantum
integrale particulare continetur, cum completum sit $y = x + ae^{x}$, denotante $e$ numerum, cuius
logarithmus est $=1$. Nisi igitur constans arbitraria $a$ euanescens ponatur, integrale semper erit
transcendens.

\origpage{39}
\textbf{§.~4.} Cum igitur euenire queat, vt aequatio differentialis integrale particulare
algebraicum admittat, etiamsi integrale completum sit transcendens, ita etiam rationes dubitandi
non desunt, quod integrale completum aequationis differentialis propositae
$\dfrac{m\,dx}{\sqrt{1-x^{4}}} = \dfrac{n\,dy}{\sqrt{1-y^{4}}}$ quantitates transcendentes inuoluat,
etiamsi pro ea integrale particulare algebraicum exhibere licuerit. Cum enim integrale completum
sit:\ednote{In the printed source the second radical of this display is set with a plus, $\sqrt{1+y^{4}}$; an apparent misprint for $\sqrt{1-y^{4}}$ --- the sign is negative in the differential equation itself and everywhere else. Reproduced here exactly as printed.}
\[ m\int \frac{dx}{\sqrt{1-x^{4}}} = n\int \frac{dy}{\sqrt{1+y^{4}}} + C \]
haec autem integralia nullo modo, neque circuli, neque hyperbolae, quadraturam in subsidium
vocando, assignari queant, minime probabile videtur, istas formulas tantopere transcendentes in
genere, ita vt constans $C$ maneat indeterminata, ad relationem algebraicam inter $x$ et $y$
reuocari posse.

\textbf{§.~5.} Notum quidem est, integrale completum huius aequationis differentialis
$\dfrac{m\,dx}{\sqrt{1-xx}} = \dfrac{n\,dy}{\sqrt{1-yy}}$ semper algebraice exhiberi posse, dummodo
proportio coefficientium $m$ et $n$ fuerit rationalis; sed quia vtriusque formulae integrale arcum
circuli indicat, ita vt integrale completum sit
$m\,\mathrm{A}\,\mathrm{sin.}\,x = n\,\mathrm{A}\,\mathrm{sin.}\,y + C$, relatio autem sinuum, qui ad
arcus proportionem rationalem inter se tenentes spectant, algebraice exprimi potest, mirum non est,
aequationem integralem completam his casibus quoque algebraice exhiberi posse. Cum autem huiusmodi
comparatio in formulis transcendentibus $\displaystyle\int \frac{dx}{\sqrt{1-x^{4}}}$ et
$\displaystyle\int \frac{dy}{\sqrt{1-y^{4}}}$ locum non habeat, seu saltem non
\origpage{40}
constet, inde reductio integralis ad quantitates algebraicas peti non poterit.

\textbf{§.~6.} Nihilo tamen minus obseruaui, si proposita fuerit huiusmodi aequatio
differentialis
\[ \frac{m\,dx}{\sqrt{1-x^{4}}} = \frac{n\,dy}{\sqrt{1-y^{4}}} \]
etiam integrale completum, quod scilicet quantitatem constantem arbitrariam inuoluat, semper
algebraice exprimi posse, dummodo ratio $m:n$ fuerit rationalis: quod mihi quidem eo magis notatu
dignum videtur, quod nulla certa methodo ad hoc integrale sum perductus, sed id potius tentando,
vel diuinando, elicui. Vnde nullum est dubium, quin methodus directa, ad idem hoc integrale
perducens, fines analyseos non mediocriter sit amplificatura; cuius propterea inuestigatio
Analystis omni studio commendanda videtur.

\textbf{§.~7.} Completum autem integrale aequationis istius differentialis, quaecunque fuerit
ratio rationalis coëfficientium $m$ et $n$, deriuare mihi licuit ex integratione completa huius
aequationis $\dfrac{dx}{\sqrt{1-x^{4}}} = \dfrac{dy}{\sqrt{1-y^{4}}}$: hac enim concessa methodum
certam indicabo, ex ea quoque integrale completum huius aequationis multo latius patentis
$\dfrac{m\,dx}{\sqrt{1-x^{4}}} = \dfrac{n\,dy}{\sqrt{1-y^{4}}}$ concludendi. Quae methodus etiam in
genere ad huiusmodi aequationum $mX\,dx = nY\,dy$ integralia inuenienda adhiberi queat, si modo
integrale completum huius $X\,dx = Y\,dy$ fuerit erutum, atque $Y$ talem significet functionem
ipsius $y$, qualis $X$ est ipsius $x$.

\origpage{41}
\textbf{§.~8.} Exordiar igitur ab hac aequatione
\[ \frac{dx}{\sqrt{1-x^{4}}} = \frac{dy}{\sqrt{1-y^{4}}} \]
cui quidem primo intuitu satisfacere perspicuum est aequationem $x = y$, quae propterea eius est
integrale particulare. Tum vero eidem aequationi quoque satisfacit iste valor algebraicus
$x = -\sqrt{\dfrac{1-yy}{1+yy}}$, cum enim sit
$dx = +\dfrac{2y\,dy}{(1+yy)\sqrt{(1-yy)(1+yy)}}$ et $\sqrt{1-x^{4}} = \dfrac{2y}{1+yy}$, erit
$\dfrac{dx}{\sqrt{1-x^{4}}} = \dfrac{dy}{\sqrt{1-y^{4}}}$. Hinc iste etiam valor, seu aequatio
$xxyy + xx + yy - 1 = 0$ est integralis particularis aequationis differentialis propositae. Vnde
integrale completum, quod constantem arbitrariam inuoluat, ita comparatum sit necesse est, vt
tribuendo huic constanti certum quendam valorem, prodeat $x = y$; sin autem eidem constanti alius
quidem valor tribuatur, vt prodeat $x = -\sqrt{\dfrac{1-yy}{1+yy}}$ seu $xxyy + xx + yy - 1 = 0$.

\subsection*{Theorema.}
\textbf{§.~9.} Dico igitur huius aequationis differentialis
\[ \frac{dx}{\sqrt{1-x^{4}}} = \frac{dy}{\sqrt{1-y^{4}}} \]
aequationem integralem completam esse:
\[ xx + yy + ccxxyy = cc + 2xy\sqrt{1-c^{4}} \]

\subsection*{Demonstratio.}
Posita enim hac aequatione, eius differentiale erit:
\[ x\,dx + y\,dy + ccxy(x\,dy + y\,dx) = (x\,dy + y\,dx)\sqrt{1-c^{4}} \]
\origpage{42}
vnde fit
\[ dx\bigl(x + ccxyy - y\sqrt{1-c^{4}}\bigr) + dy\bigl(y + ccxxy - x\sqrt{1-c^{4}}\bigr) = 0 \]
Ex eadem vero aequatione resoluta colligitur:
\[ y = \frac{x\sqrt{1-c^{4}} + c\sqrt{1-x^{4}}}{1+ccxx} \quad\text{et}\quad
   x = \frac{y\sqrt{1-c^{4}} - c\sqrt{1-y^{4}}}{1+ccyy} \]
Si enim ibi radicali $\sqrt{1-x^{4}}$ tribuitur signum $+$, hic radicali $\sqrt{1-y^{4}}$ signum $-$
tribui debet, vt posito $x = 0$, vtrinque idem valor prodeat $y = c$. Erit ergo
\[ x + ccxyy - y\sqrt{1-c^{4}} = -c\sqrt{1-y^{4}} \]
\[ y + ccxxy - x\sqrt{1-c^{4}} = c\sqrt{1-x^{4}} \]
quibus valoribus in aequatione differentiali substitutis, prodit
\[ -c\,dx\sqrt{1-y^{4}} + c\,dy\sqrt{1-x^{4}} = 0, \]
siue
\[ \frac{dx}{\sqrt{1-x^{4}}} = \frac{dy}{\sqrt{1-y^{4}}}. \]
Huius ergo aequationis differentialis integrale est:
\[ xx + yy + ccxxyy = cc + 2xy\sqrt{1-c^{4}} \]
et quia constantem $c$ ab arbitrio nostro pendentem continet, erit simul integrale completum.
Q.~E.~D.

\textbf{§.~10.} Si igitur habeatur haec aequatio
$\dfrac{dx}{\sqrt{1-x^{4}}} = \dfrac{dy}{\sqrt{1-y^{4}}}$ valor integralis completus ipsius $x$ est:
\[ x = \frac{y\sqrt{1-c^{4}} + c\sqrt{1-y^{4}}}{1+ccyy} \]
vnde si constans arbitraria $c$ euanescat fit $x = y$; sin autem ponatur $c = 1$, habemus
$x = \pm\dfrac{\sqrt{1-y^{4}}}{1+yy} = \sqrt{\dfrac{1-yy}{1+yy}}$ qui sunt ambo illi valores
particulares iam supra exhibiti.

\origpage{43}
Hinc eruuntur alii valores particulares prae caeteris simpliciores, sed qui ad imaginaria
deuoluuntur. Ita posito $c = \infty$ fit $x = \dfrac{\sqrt{-1}}{y}$; et posito $cc = -1$; fit
$x = \sqrt{\dfrac{yy+1}{yy-1}}$ qui itidem aequationi propositae satisfaciunt.

\textbf{§.~11.} Quo autem ratio huius integralis clarius perspiciatur, concipiatur curua $AM$
cuius haec sit indoles, vt posita abscissa $AP = u$, sit arcus ei respondens
$AM = \displaystyle\int \frac{du}{\sqrt{1-u^{4}}}$.
\rmfigure{figures/fig-1.png}{Fig.~1.}{Fig. 1 (Tab. I): the curve AM with abscissa AP; the arc from A measures the elliptic integral treated in the paper. Cropped from the original engraved plate, Novi Commentarii vol. VI (1761).}
Deinde eadem curua denuo descripta, capiatur abscissa
$ap = x$, erit arcus $am = \displaystyle\int \frac{dx}{\sqrt{1-x^{4}}}$. Sumto igitur
\[ x = \frac{u\sqrt{1-c^{4}} + c\sqrt{1-u^{4}}}{1+ccuu} \]
fiet $\dfrac{dx}{\sqrt{1-x^{4}}} = \dfrac{du}{\sqrt{1-u^{4}}}$; ideoque arc.~$am = $ arc.~$AM + $
Const. Pro constantis autem huius determinatione, posito $u = 0$ quo casu arcus $AM$ euanescit, fit
$x = c$. Quare si capiatur abscissa $ab = c$, cui arcus $ad$ respondeat, erit arcus $dm = $ arcui
$AM$.

\textbf{§.~12.} Ope huius ergo integrationis completa aequationis
$\dfrac{dx}{\sqrt{1-x^{4}}} = \dfrac{du}{\sqrt{1-u^{4}}}$, in curua proposita arcui cuicunque $AM$,
qui abscissae $AP = u$ respondet, arcus aequalis $dm$, qui a dato puncto $d$ incipiat, abscindi
poterit. Posita enim abscissa dato puncto $d$ respondente $ab = c$; si capiatur abscissa
$ap = x = \dfrac{c\sqrt{1-u^{4}} + u\sqrt{1-c^{4}}}{1+ccuu}$ erit arcus $dm$ arcui $AM$ aequalis.
\rmfigure{figures/fig-2.png}{Fig.~2.}{Fig. 2 (Tab. I): a second copy of the curve, with the abscissae ap and a-pi (points a, pi, b, p, d, m) by which an arc dm or d-mu equal to a given arc AM is cut off from the fixed point d. Cropped from the original engraved plate, Novi Commentarii vol. VI (1761).}
Simili autem modo cum $\sqrt{1-c^{4}}$ negatiuum statui liceat, si capiatur abscissa
\[ a\pi = \frac{c\sqrt{1-u^{4}} - u\sqrt{1-c^{4}}}{1+ccuu} \]
\origpage{44}
erit itidem arcus $d\mu$ arcui $AM$ aequalis: sicque in hac curua a dato quouis puncto $d$ vtrinque
abscindi potest arcus $dm$ et $d\mu$, qui arcui $AM$ sint aequales.

\textbf{§.~13.} Hinc ergo patet, si arcus $ad$ aequalis capiatur arcui $AM$, seu $c = u$, fore
arcum $am$ duplum arcus $AM$. Hinc si statuatur $ap = x = \dfrac{2u\sqrt{1-u^{4}}}{1+u^{4}}$,
prodibit arcus $am = 2$~arc.~$AM$. Simili modo si capiatur arcus $ad = 2AM$, seu
$c = \dfrac{2u\sqrt{1-u^{4}}}{1+u^{4}}$, statuaturque
$x = \dfrac{c\sqrt{1-u^{4}} + u\sqrt{1-c^{4}}}{1+ccuu}$ obtinebitur arcus $am = 3$~arc.~$AM$. Ac si
iste valor ipsius $x$ denuo pro $c$ substituatur, vt fit $ad = 3AM$ iterumque statuatur
$x = \dfrac{c\sqrt{1-u^{4}} + u\sqrt{1-c^{4}}}{1+ccuu}$, nascetur arcus $am$ quadruplus arcus $AM$;
atque ita porro successiue quaecunque multipla arcus $AM$ geometrice assignari poterunt.

\textbf{§.~14.} Sit arcus $ad = n \cdot AM$ et $ab = z$; ita vt sit
$\displaystyle\int \frac{dz}{\sqrt{1-z^{4}}} = n\int \frac{du}{\sqrt{1-u^{4}}}$; atque ex his patet si capiatur
$x = \dfrac{z\sqrt{1-u^{4}} + u\sqrt{1-z^{4}}}{1+uuzz}$ fore
$\displaystyle\int \frac{dx}{\sqrt{1-x^{4}}} = (n+1)\int \frac{du}{\sqrt{1-u^{4}}}$; sin autem ponatur
$x = \dfrac{z\sqrt{1-u^{4}} - u\sqrt{1-z^{4}}}{1+uuzz}$, tum futurum esse
$\displaystyle\int \frac{dx}{\sqrt{1-x^{4}}} = (n-1)\int \frac{du}{\sqrt{1-u^{4}}}$. Si igitur haec aequatio
$\dfrac{dz}{\sqrt{1-z^{4}}} = \dfrac{n\,du}{\sqrt{1-u^{4}}}$ fuerit integrata, debitusque valor pro
$z$ inde erutus, etiam integrari poterit haec aequatio
$\dfrac{dx}{\sqrt{1-x^{4}}} = \dfrac{(n\pm1)\,du}{\sqrt{1-u^{4}}}$, quippe cuius integrale erit
$x = \dfrac{z\sqrt{1-u^{4}} \pm u\sqrt{1-z^{4}}}{1+uuzz}$. Ac si pro $z$ assumtus fuerit eius valor
completus, qui scilicet constantem arbitrariam inuoluat, etiam pro $x$ prodibit eius valor
completus.

\origpage{45}
\textbf{§.~15.} Hinc igitur perspicuum est, quomodo aequatio integralis completa inueniri
debeat, quae conueniat huic aequationi differentiali
$\dfrac{dx}{\sqrt{1-x^{4}}} = \dfrac{n\,du}{\sqrt{1-u^{4}}}$; quoties $n$ fuerit numerus integer.
Simili autem modo assignari poterit $y$, vt fit
$\dfrac{dy}{\sqrt{1-y^{4}}} = \dfrac{m\,du}{\sqrt{1-u^{4}}}$, vnde si eliminando $u$, aequatio inter
$x$ et $y$ quaeratur, ea erit integralis huius aequationis
$\dfrac{m\,dx}{\sqrt{1-x^{4}}} = \dfrac{n\,dy}{\sqrt{1-y^{4}}}$, quicunque numeri rationales pro $m$
et $n$ substituantur: atque vt hoc integrale prodeat completum, sufficit pro altera tantum
variabilium $x$ et $y$ valorem completum per $u$ determinasse, cum hinc iam noua constans arbitraria
in calculum introducatur.

\textbf{§.~16.} Methodus, qua hic in Theorematis demonstratione sum vsus, etsi non ex rei
natura est petita, sed indirecte ad id, quod propositum erat, perduxit, tamen multo latius patet:
simili enim modo colligitur, huius aequationis differentialis
\[ \frac{dx}{\sqrt{1+mxx+nx^{4}}} = \frac{dy}{\sqrt{1+myy+ny^{4}}} \]
integrale completum esse:
\[ 0 = cc - xx - yy + nccxxyy + 2xy\sqrt{1+mcc+nc^{4}} \]
Vnde idem, quod ante, ratiocinium adhibendo, integrale quoque completum obtinebitur huius
aequationis
\[ \frac{\mu\,dx}{\sqrt{1+mxx+nx^{4}}} = \frac{\nu\,dy}{\sqrt{1+myy+ny^{4}}} \]
siquidem litteris $\mu$ et $\nu$ numeri integri designentur.

\textbf{§.~17.} Inuestigatio autem huius integrationis ita se habet: Fingatur primo pro
arbitrio relatio inter variabiles $x$ et $y$ hac aequatione contenta:
\origpage{46}
\[ \alpha xx + \alpha yy = 2\beta xy + \gamma xxyy + \delta \tag{1} \]
quae differentiata dat:
\[ \alpha x\,dx + \alpha y\,dy = \beta x\,dy + \beta y\,dx + \gamma xyy\,dx + \gamma xxy\,dy \]
vnde conficitur
\[ dx(\alpha x - \beta y - \gamma xyy) + dy(\alpha y - \beta x - \gamma xxy) = 0 \tag{2} \]
Deinde ex aequatione (1) eliciantur valores vtriusque variabilis:
\[ x = \frac{\beta y + \sqrt{\alpha\delta + (\beta\beta - \alpha\alpha - \gamma\delta)yy + \alpha\gamma y^{4}}}{\alpha - \gamma yy} \]
\[ y = \frac{\beta x - \sqrt{\alpha\delta + (\beta\beta - \alpha\alpha - \gamma\delta)xx + \alpha\gamma x^{4}}}{\alpha - \gamma xx} \]
Atque hinc obtinemus:
\[ \alpha x - \beta y - \gamma xyy = \sqrt{\alpha\delta + (\beta\beta - \alpha\alpha - \gamma\delta)yy + \alpha\gamma y^{4}} \tag{3} \]
\[ \alpha y - \beta x - \gamma xxy = -\sqrt{\alpha\delta + (\beta\beta - \alpha\alpha - \gamma\delta)xx + \alpha\gamma x^{4}} \tag{4} \]
qui valores in aequatione (2) substituti praebebunt
\[ \frac{dx}{\sqrt{\alpha\delta + (\beta\beta - \alpha\alpha - \gamma\delta)xx + \alpha\gamma x^{4}}}
   = \frac{dy}{\sqrt{\alpha\delta + (\beta\beta - \alpha\alpha - \gamma\delta)yy + \alpha\gamma y^{4}}} \tag{5} \]
cuius ergo aequationis integrale est aequatio (1).

\textbf{§.~18.} Quo istas formas simpliciores reddamus, ponamus $\alpha\delta = A$;
$\beta\beta - \alpha\alpha - \gamma\delta = C$; $\alpha\gamma = E$ eritque $\delta = \dfrac{A}{\alpha}$;
$\gamma = \dfrac{E}{\alpha}$ et $\beta = \sqrt{C + \alpha\alpha + \dfrac{AE}{\alpha\alpha}}$. Quare huius
aequationis differentialis
\[ \frac{dx}{\sqrt{A + Cxx + Ex^{4}}} = \frac{dy}{\sqrt{A + Cyy + Ey^{4}}} \tag{6} \]
aequatio integralis est haec:
\[ \alpha(xx + yy) = \frac{A}{\alpha} + \frac{E}{\alpha}xxyy + 2xy\sqrt{C + \alpha\alpha + \frac{AE}{\alpha\alpha}} \tag{7} \]
quae simul est integralis completa:

\textbf{§.~19.} Vel ponamus $A = f\alpha\alpha$: $C = g\alpha\alpha$ et $E = b\alpha\alpha$, vt
habeamus hanc aequationem differentialem
\[ \frac{dx}{\sqrt{f + gxx + bx^{4}}} = \frac{dy}{\sqrt{f + gyy + by^{4}}} \]
\origpage{47}
cuius propterea aequatio integralis completa erit:
\[ xx + yy = f + bxxyy + 2xy\sqrt{1 + g + fb} \]
quae etsi nouam constantem inuoluere non videtur, tamen est completa, cum in differentiali tantum
ratio quantitatum $f$, $g$, et $b$ spectatur, ita vt pro $f$, $g$, et $b$ scribere liceat $fcc$,
$gcc$ et $bcc$, vnde aequatio integralis manifesto completa prodit:
\[ xx + yy = fcc + bccxxyy + 2xy\sqrt{1 + gcc + fbc^{4}} \]
vel $f(xx + yy) = fee + beexxyy + 2xy\sqrt{f(f + gee + be^{4})}$ posito $cc = \dfrac{ee}{f}$.

\textbf{§.~20.} Quodsi ergo proposita sit haec aequatio differentialis
\[ \frac{dx}{\sqrt{f + gxx + bx^{4}}} = \frac{dy}{\sqrt{f + gyy + by^{4}}} \]
valor ipsius $y$ per functionem algebraicam ipsius $x$ exprimi poterit, ita vt sit:
\[ y = \frac{x\sqrt{1 + gcc + fbc^{4}} \pm c\sqrt{f + gxx + bx^{4}}}{1 - bccxx} \]
vel $y = \dfrac{x\sqrt{f(f + gee + be^{4})} \pm e\sqrt{f(f + gxx + bx^{4})}}{f - beexx}$.

Quodsi ergo sit $g = 0$, vt habeatur haec aequatio differentialis
\[ \frac{dx}{\sqrt{f + bx^{4}}} = \frac{dy}{\sqrt{f + by^{4}}} \]
valor integralis completus ipsius $y$ erit
\[ y = \frac{x\sqrt{f(f + be^{4})} \pm e\sqrt{f(f + bx^{4})}}{f - beexx} \]
vnde constantem $e$ pro lubitu determinando, innumeri valores particulares pro $y$ deduci possunt.

\origpage{48}
\textbf{§.~21.} Methodi autem, qua supra vsus sum, beneficio etiam huius aequationis
\[ \frac{m\,dx}{\sqrt{f + gxx + bx^{4}}} = \frac{n\,dy}{\sqrt{f + gyy + by^{4}}} \]
si modo $m$ et $n$ sint numeri rationales, integrale completum, atque id quidem algebraice,
exhiberi poterit.

\textbf{§.~22.} Quemadmodum in aequatione supra assumta, variabiles $x$ et $y$ inter se
permutabiles sunt constitutae, vt ambae formulae inter se similes euaderent, ita omissa hac
limitatione ad formularum differentialium disparium comparationem perueniemus. Ponamus ergo:
\[ \alpha xx + \beta yy = 2\gamma xy + \delta xxyy + \varepsilon \tag{1} \]
vnde fit
\[ x = \frac{\gamma y + \sqrt{\alpha\varepsilon + (\gamma\gamma - \delta\varepsilon - \alpha\beta)yy + \beta\delta y^{4}}}{\alpha - \delta yy} \]
et $y = \dfrac{\gamma x - \sqrt{\beta\varepsilon + (\gamma\gamma - \delta\varepsilon - \alpha\beta)xx + \alpha\delta x^{4}}}{\beta - \delta xx}$
hincque
\[ \alpha x - \gamma y - \delta xyy = \sqrt{\alpha\varepsilon + (\gamma\gamma - \delta\varepsilon - \alpha\beta)yy + \beta\delta y^{4}} \tag{2} \]
\[ \beta y - \gamma x - \delta xxy = -\sqrt{\beta\varepsilon + (\gamma\gamma - \delta\varepsilon - \alpha\beta)xx + \alpha\delta x^{4}} \tag{3} \]
at aequatio (1) differentiata dat:
\[ dx(\alpha x - \gamma y - \delta xyy) + dy(\beta y - \gamma x - \delta xxy) = 0 \]
vnde conficitur haec aequatio differentialis:
\[ \frac{dx}{\sqrt{\beta\varepsilon + (\gamma\gamma - \delta\varepsilon - \alpha\beta)xx + \alpha\delta x^{4}}}
   = \frac{dy}{\sqrt{\alpha\varepsilon + (\gamma\gamma - \delta\varepsilon - \alpha\beta)yy + \beta\delta y^{4}}} \]
cuius propterea integralis est aequatio assumta.

\textbf{§.~23.} Verum haec disparitas facile tollitur, loco $y$ ponendo $z\sqrt{\dfrac{\alpha}{\beta}}$,
cuius rei ratio statim ex aequatione assumta potuisset esse manifesta. Sed alia patet via ad
formulas dispares perueniendi, cuius hic exemplum tradidisse sufficiat. Assumatur aequatio:
$x^{4} + 2axxyy + 2bxx$
\origpage{49}
${}= c$, cuius differentiale est $dx(x^{3} + axyy + bx) + axxy\,dy = 0$, seu
\[ \frac{dx}{xy} = \frac{-a\,dy}{xx + ayy + b} \]
Iam ex aequatione assumta primo determinetur $xy$ per $x$ sicque fiet
$xy = \sqrt{\dfrac{c - 2bxx - x^{4}}{2a}}$; tum vero $xx + ayy + b$ per $y$, at ob
$(xx + ayy + b)^{2} = c + (ayy + b)^{2}$, erit
\[ xx + ayy + b = \sqrt{c + (ayy + b)^{2}} \]
Quocirca habebitur aequatio differentialis ista
\[ \frac{dx\sqrt{2a}}{\sqrt{c - 2bxx - x^{4}}} = \frac{a\,dy}{\sqrt{c + bb + 2abyy + aay^{4}}} \]
cuius propterea integralis est assumta seu $y = \dfrac{\sqrt{c - 2bxx - x^{4}}}{x\sqrt{2a}}$.

\textbf{§.~24.} Etsi hoc integrale non est completum, tamen ex superioribus facile completum
reddetur. Ponatur enim:
\[ \frac{a\,dy}{\sqrt{c + bb + 2abyy + aay^{4}}} = \frac{a\,dz}{\sqrt{c + bb + 2abzz + aaz^{4}}} \]
ob $f = c + bb$; $g = 2ab$; $b = aa$, erit
\[ y = \frac{z\sqrt{(c+bb)(c + bb + 2abee + aae^{4})} \pm e\sqrt{(c+bb)(c + bb + 2abzz + aaz^{4})}}{c + bb - aaeezz} \]
hic ergo valor aequalis statuatur ipsi $\dfrac{\sqrt{c - 2bxx - x^{4}}}{x\sqrt{2a}}$, et aequatio hinc
inter $x$ et $z$ resultans integralis erit completa huius aequationi differentialis
\[ \frac{dx\sqrt{2a}}{\sqrt{c - 2bxx - x^{4}}} = \frac{a\,dz}{\sqrt{c + bb + 2abzz + aaz^{4}}} \]
Quin etiam ex allatis patet, si haec bina membra insuper per numeros rationales quoscunque
multiplicentur, quemadmodum integrale completum inueniri oporteat.

\origpage{50}
\textbf{§.~25.} Verum missa membrorum disparitate formationem parium membrorum generalius
concipiamus, ponatur ergo:
\[ 0 = \alpha + 2\beta(x+y) + \gamma(xx+yy) + 2\delta xy + 2\varepsilon xy(x+y) + \zeta xxyy \tag{1} \]
vnde differentiando obtinetur:
\[ dx(\beta + \gamma x + \delta y + 2\varepsilon xy + \varepsilon yy + \zeta xyy)
   + dy(\beta + \gamma y + \delta x + 2\varepsilon xy + \varepsilon xx + \zeta xxy) = 0 \]
ideoque
\[ \frac{dy}{\beta + \gamma x + \delta y + 2\varepsilon xy + \varepsilon yy + \zeta xyy}
   = \frac{-dx}{\beta + \gamma y + \delta x + 2\varepsilon xy + \varepsilon xx + \zeta xxy} \tag{2} \]
Ex resolutione autem aequationis assumtae elicitur.
\[ y = \frac{-\beta - \delta x - \varepsilon xx + \sqrt{\beta\beta - \alpha\gamma + 2(\beta\delta - \alpha\varepsilon - \beta\gamma)x + (\delta\delta - \gamma\gamma - \alpha\zeta - 2\beta\varepsilon)xx + 2(\delta\varepsilon - \beta\zeta - \gamma\varepsilon)x^{3} + (\varepsilon\varepsilon - \gamma\zeta)x^{4}}}{\gamma + 2\varepsilon x + \zeta xx} \]
Ponatur breuitatis gratia
\[ \beta\beta - \alpha\gamma = A; \quad \beta\delta - \alpha\varepsilon - \beta\gamma = B; \quad
   \delta\delta - \gamma\gamma - \alpha\zeta - 2\beta\varepsilon = C; \quad
   \varepsilon\varepsilon - \gamma\zeta = E; \quad \delta\varepsilon - \beta\zeta - \gamma\varepsilon = D; \]
eritque
\[ \beta + \delta x + \varepsilon xx + \gamma y + 2\varepsilon xy + \zeta xxy = \pm\sqrt{A + 2Bx + Cxx + 2Dx^{3} + Ex^{4}} \]
\[ \beta + \delta y + \varepsilon yy + \gamma x + 2\varepsilon xy + \zeta xyy = \mp\sqrt{A + 2By + Cyy + 2Dy^{3} + Ey^{4}} \]

\textbf{§.~26.} Hinc itaque concludimus huius aequationis differentialis:
\[ \frac{dx}{\sqrt{A + 2Bx + Cxx + 2Dx^{3} + Ex^{4}}} = \frac{dy}{\sqrt{A + 2By + Cyy + 2Dy^{3} + Ey^{4}}} \]
aequationem integralem eamque completam esse
\[ 0 = \alpha + 2\beta(x+y) + \gamma(xx+yy) + 2\delta xy + 2\varepsilon xy(x+y) + \zeta xxyy \]
\origpage{51}
adhibita scilicet superiori horum coëfficientium determinatione. Primum autem definiatur $\beta$
vel $\varepsilon$ ex hac aequatione
\[ \frac{BB(\varepsilon\varepsilon - E) - DD(\beta\beta - A)}{A\varepsilon\varepsilon - E\beta\beta}
   + \frac{2AD\varepsilon - 2BE\beta}{B\varepsilon - D\beta} = C \]
tum vero erit:
\[ \gamma = \frac{A\varepsilon\varepsilon - E\beta\beta}{B\varepsilon - D\beta}; \quad
   \alpha = \frac{\beta\beta - A}{\gamma}; \quad \zeta = \frac{\varepsilon\varepsilon - E}{\gamma} \quad\text{et} \]
\[ \delta = \frac{B\beta(\varepsilon\varepsilon - E) - D\varepsilon(\beta\beta - A)}{A\varepsilon\varepsilon - E\beta\beta} + \gamma
   \quad\text{seu}\quad \delta = \gamma + \frac{B + \alpha\varepsilon}{\beta} \]

\textbf{§.~27.} Hinc ergo perspicuum est etiam hanc aequationem differentialem:
\[ \frac{dx}{\sqrt{A + 2Dx^{3}}} = \frac{dy}{\sqrt{A + 2Dy^{3}}} \]
integrari posse: nam ob $B = 0$, $C = 0$ et $E = 0$ erit
\[ -\frac{DD(\beta\beta - A)}{A\varepsilon\varepsilon} - \frac{2A\varepsilon}{\beta} = 0 \quad\text{seu}\quad
   \varepsilon = \sqrt[3]{\frac{DD}{2AA}\beta(A - \beta\beta)} \]
at hinc valores nimis prodeunt complicati. Facilius negotium absoluetur, resoluendo valores
litterarum euanescentium $B$, $C$ et $E$; nam $E = 0$ dat: $\zeta = \dfrac{\varepsilon\varepsilon}{\gamma}$;
tum $B = 0$ dat: $\delta = \gamma + \dfrac{\alpha\varepsilon}{\beta}$; atque $C = 0$ dat
$\delta\delta - \gamma\gamma = \alpha\zeta + 2\beta\varepsilon = \dfrac{\alpha\varepsilon\varepsilon}{\gamma} + 2\beta\varepsilon
= \dfrac{\alpha^{2}\varepsilon\varepsilon}{\beta\beta} + \dfrac{2\alpha\gamma\varepsilon}{\beta}$ cuius factores sunt
$\beta\beta = \alpha\gamma$ et $\alpha\varepsilon\varepsilon + 2\beta\gamma\varepsilon = 0$. At si esset
$\beta\beta = \alpha\gamma$ foret $A = 0$, sin autem esset $\varepsilon = 0$ foret et $\zeta = 0$ et
$D = 0$, contra scopum. Fieri ergo oportet $\alpha\varepsilon = -2\beta\gamma$; vnde fiet
$\alpha = -\dfrac{2\beta\gamma}{\varepsilon}$; $\delta = -\gamma$; et $\zeta = \dfrac{\varepsilon\varepsilon}{\gamma}$.
Denique fieri debet $\beta\beta + \dfrac{2\beta\gamma\gamma}{\varepsilon} = A$ et
$-2\gamma\varepsilon - \dfrac{\beta\varepsilon\varepsilon}{\gamma} = D$. Inde fit
$\varepsilon = \dfrac{2\beta\gamma\gamma}{A - \beta\beta}$; et ob
$\dfrac{\gamma D}{\varepsilon} = -(2\gamma\gamma + \beta\varepsilon)$ et $2\gamma\gamma + \beta\varepsilon = \dfrac{A\varepsilon}{\beta}$,
erit $\dfrac{\gamma D}{\varepsilon} = -\dfrac{A\varepsilon}{\beta}$; ideoque $\varepsilon\varepsilon = -\dfrac{\beta\gamma D}{A}$. Ergo
\[ \frac{4\beta\gamma^{3}}{(A - \beta\beta)^{2}} + \frac{D}{A} = 0. \]

\origpage{52}
\textbf{§.~28.} Cum autem tantum ratio litterarum $A$ et $D$ in censum veniat, aequatio vltima
valori absoluto ipsius $A$ inueniendo inseruit, quem autem nosse non est opus. Manebunt ergo
litterae $\gamma$ et $\beta$ indeterminatae. Ponatur ergo $\gamma = -Ac$ et $\beta = Dc$, erit
$\varepsilon\varepsilon = DDcc$, seu $\varepsilon = Dc$, hincque $\delta = Ac$;
$\zeta = -\dfrac{DDc}{A}$; et $\alpha = 2Ac$. Quare huius aequationis differentialis:
\[ \frac{dx}{\sqrt{A + 2Dx^{3}}} = \frac{dy}{\sqrt{A + 2Dy^{3}}} \]
integrale est.
\[ 0 = 2A + 2D(x+y) - A(xx+yy) + 2Axy + 2Dxy(x+y) - \frac{DD}{A}xxyy \]
Hoc autem integrale non est completum, tale autem reddetur ponendo $\gamma = -A$ et $\beta = Dcc$,
vnde fit $\varepsilon\varepsilon = DDcc$ et $\varepsilon = Dc$; porro erit $\delta = A$;
$\zeta = -\dfrac{DDcc}{A}$; $\alpha = 2Ac$; ita vt integrale completum sit:
\[ 0 = 2Ac + 2Dcc(x+y) - A(xx+yy) + 2Axy + 2Dcxy(x+y) - \frac{DDcc}{A}xxyy \]
vbi $c$ est constans ab arbitrio pendens, vnde fit:
\[ y = \frac{Dcc + Ax + Dcxx + \sqrt{c\bigl(2A + \frac{DD}{A}c^{3}\bigr)(A + 2Dx^{3})}}{A - 2Dcx + \frac{DDcc}{A}xx} \]

\textbf{§.~29.} Hic casus notari meretur, quo $A = 1$ et $D = \frac{1}{2}$, vt habeatur haec
aequatio differentialis
\[ \frac{dx}{\sqrt{1 + x^{3}}} = \frac{dy}{\sqrt{1 + y^{3}}} \]
vbi ad fractiones tollendas loco $c$ scribatur $2c$ eritque integrale completum:
\[ 0 = 4c + 4cc(x+y) - xx - yy + 2xy + 2cxy(x+y) - ccxxyy \]
seu $y = \dfrac{2cc + x + cxx + 2\sqrt{c(1+c^{3})(1+x^{3})}}{1 - 2cx + ccxx}$

\origpage{53}
Integralia ergo particularia erunt

I.~si $c = 0$; $y = x$;

II.~si $c = \infty$; $y = \dfrac{2 \pm 2\sqrt{1+x^{3}}}{xx}$;

III.~si $c = -1$; $y = \dfrac{2 + x - xx}{1 + 2x + xx} = \dfrac{2-x}{1+x}$.

\textbf{§.~30.} Ex eodem principio si in §.~29. loco litterarum $A$, $B$, $C$, $D$, $E$,
eaedem per quantitatem quampiam $p$ multiplicentur, nihilo minus aequatio differentialis erit
\[ \frac{dx}{\sqrt{A + 2Bx + Cxx + 2Dx^{3} + Ex^{4}}} = \frac{dy}{\sqrt{A + 2By + Cyy + 2Dy^{3} + Ey^{4}}} \]
inuenieturque
\[ p = \frac{BB\varepsilon\varepsilon - DD\beta\beta}{BBE - ADD}
   + 2\frac{(AD\varepsilon - BE\beta)(A\varepsilon\varepsilon - E\beta\beta)}{(B\varepsilon - D\beta)(BBE - ADD)}
   - \frac{C(A\varepsilon\varepsilon - E\beta\beta)}{BBE - ADD} \]
tum erit $\gamma = \dfrac{A\varepsilon\varepsilon - E\beta\beta}{B\varepsilon - D\beta}$;
$\alpha = \dfrac{\beta\beta - Ap}{\gamma}$; $\zeta = \dfrac{\varepsilon\varepsilon - Ep}{\gamma}$ atque
$\delta = \gamma + \dfrac{\alpha\varepsilon + Bp}{\beta}$: ita vt litterae $\beta$ et $\varepsilon$ maneant
indeterminatae, fietque propterea aequatio integralis completa:
\[ 0 = \alpha + 2\beta(x+y) + \gamma(xx+yy) + 2\delta xy + 2\varepsilon xy(x+y) + \zeta xxyy \]
vnde fit:
\[ y = \frac{-\beta - \delta x - \varepsilon xx + \sqrt{p(A + 2Bx + Cxx + 2Dx^{3} + Ex^{4})}}{\gamma + 2\varepsilon x + \zeta xx} \]

\textbf{§.~31.} Notandum denique est, non solum hanc aequationem differentialem, cuius
integrale completum modo exhibui, sed etiam hanc multo latius patentem
\[ \frac{m\,dx}{\sqrt{A + 2Bx + Cxx + 2Dx^{3} + Ex^{4}}} = \frac{n\,dy}{\sqrt{A + 2By + Cyy + 2Dy^{3} + Ey^{4}}} \]
semper algebraice et quidem complete integrari posse, dummodo coëfficientium $m$ et $n$ ratio
fuerit rationalis: haec enim integratio simili modo instituitur, quo supra vsus sum ad aequationem,
quae mihi hic praecipue erat proposita, integrandam. Methodus autem, cuius
\origpage{54}
hic specimina attuli, ita mihi videtur comparata, vt indolem eius diligentius excolendo, ad
insignes vsus apta reddi queat, vnde haud contemnenda commoda in Analysin sint redundatura.

% [reading] The 9th coefficient (constant term) prints as an italic x-like glyph; read as
% varkappa (\varkappa) — the Greek letter after theta in Euler's sequence, NOT the variable x.
% Confirmed correct at review (DvdL, 2026-07-25).
% [notation] The second set of coefficients is printed in Fraktur (distinct variables from roman
% A..E); preserved with \mathfrak per HOUSESTYLE R13, not normalized to roman.
\textbf{§.~32.} Hic autem obseruo, formulam §.~28 assumtam latius extendendo, eiusmodi
differentialia inter se comparari posse, quae sint disparia, atque adeo exemplum disparitatis
§.~26. allatum hoc modo obtineri posse; ita vt omnia, quae hactenus sunt tradita, in hac generali
inuestigatione contineantur. Fingatur scilicet haec aequatio integralis:
\[ \alpha xxyy + 2\beta xxy + 2\gamma xyy + \delta xx + \varepsilon yy + 2\zeta xy + 2\eta x + 2\theta y + \varkappa = 0 \tag{1} \]
ex qua fit
\[ y = \frac{-\beta xx - \zeta x - \theta + \sqrt{(\beta xx + \zeta x + \theta)^{2} - (\alpha xx + 2\gamma x + \varepsilon)(\delta xx + 2\eta x + \varkappa)}}{\alpha xx + 2\gamma x + \varepsilon} \tag{2} \]
\[ x = \frac{-\gamma yy - \zeta y - \eta - \sqrt{(\gamma yy + \zeta y + \eta)^{2} - (\alpha yy + 2\beta y + \delta)(\varepsilon yy + 2\theta y + \varkappa)}}{\alpha yy + 2\beta y + \delta} \tag{3} \]
Ponatur iam breuitatis gratia:
\[ \begin{aligned}
App &= \beta\beta - \alpha\delta &\qquad \mathfrak{A}qq &= \gamma\gamma - \alpha\varepsilon \\
2Bpp &= 2\beta\zeta - 2\alpha\eta - 2\gamma\delta &\qquad 2\mathfrak{B}qq &= 2\gamma\zeta - 2\alpha\theta - 2\beta\varepsilon \\
Cpp &= \zeta\zeta + 2\beta\theta - \alpha\varkappa - \delta\varepsilon - 4\gamma\eta &\qquad \mathfrak{C}qq &= \zeta\zeta + 2\gamma\eta - \alpha\varkappa - \delta\varepsilon - 4\beta\theta \\
2Dpp &= 2\zeta\theta - 2\gamma\varkappa - 2\varepsilon\eta &\qquad 2\mathfrak{D}qq &= 2\zeta\eta - 2\beta\varkappa - 2\delta\theta \\
Epp &= \theta\theta - \varepsilon\varkappa &\qquad \mathfrak{E}qq &= \eta\eta - \delta\varkappa
\end{aligned} \]
eritque:
\[ p\sqrt{Ax^{4} + 2Bx^{3} + Cxx + 2Dx + E} = \alpha xxy + 2\gamma xy + \varepsilon y + \beta xx + \zeta x + \theta \tag{4} \]
\[ -q\sqrt{\mathfrak{A}y^{4} + 2\mathfrak{B}y^{3} + \mathfrak{C}yy + 2\mathfrak{D}y + \mathfrak{E}} = \alpha xyy + 2\beta xy + \delta x + \gamma yy + \zeta y + \eta \tag{5} \]

\origpage{55}
\textbf{§.~33.} At si aequatio integralis assumta differentietur, fiet
\[ dx(\alpha xyy + 2\beta xy + \gamma yy + \delta x + \zeta y + \eta)
   + dy(\alpha xxy + \beta xx + 2\gamma xy + \varepsilon y + \zeta x + \theta) = 0 \tag{6} \]
vnde si istorum factorum valores (4) et (5) reperti substituantur, orietur ista aequatio
differentialis:
\[ \frac{q\,dx}{\sqrt{Ax^{4} + 2Bx^{3} + Cxx + 2Dx + E}}
   = \frac{p\,dy}{\sqrt{\mathfrak{A}y^{4} + 2\mathfrak{B}y^{3} + \mathfrak{C}yy + 2\mathfrak{D}y + \mathfrak{E}}} \tag{7} \]
cuius propterea integralis est aequatio assumta (1). Cum autem supra habeantur 10 aequationes,
coëfficientium autem $\alpha$, $\beta$, $\gamma$, $\delta$, etc. numerus sit 9, quorum vnus pro
lubitu assumi potest, octo remanebunt litterae determinandae. Porro autem insuper definiendae
accedunt binae litterae $p$ et $q$, ita vt nunc decem quantitates adsint incognitae; ex quo
coëfficientes vtriusque formulae $A$, $B$, $C$, $D$, $E$ et $\mathfrak{A}$, $\mathfrak{B}$,
$\mathfrak{C}$, $\mathfrak{D}$, $\mathfrak{E}$ videntur pro lubitu assumi posse. Verum perspicuum
est, cum alteri iam fuerint ad libitum assumti, alteros non omnino ab arbitrio nostro pendere,
alias enim quaeuis formula ad algebraicam reduci posset.

\textbf{§.~34.} Hinc autem aliae datae formulae transmutationes non inelegantes obtineri
possunt, si loco $y$ alii valores substituantur. Veluti si ponatur $\mathfrak{C} = 0$,\ednote{The text sets $\mathfrak C = 0$, but the gloss $\eta\eta = \delta\varkappa$ is in fact the condition $\mathfrak E = 0$ (since $\mathfrak E qq = \eta\eta - \delta\varkappa$); equation (8) that follows holds only when $\mathfrak E = 0$. An apparent misprint for $\mathfrak E = 0$, reproduced here as printed.} seu
$\eta\eta = \delta\varkappa$, statuaturque $y = zz$ sequens prodibit aequatio differentialis.
\[ \frac{q\,dx}{\sqrt{Ax^{4} + 2Bx^{3} + Cxx + 2Dx + E}}
   = \frac{2p\,dz}{\sqrt{\mathfrak{A}z^{6} + 2\mathfrak{B}z^{4} + \mathfrak{C}z^{2} + 2\mathfrak{D}}} \tag{8} \]
cuius propterea integralis est aequatio assumta, si ponatur $y = zz$, statuaturque
$\eta\eta = \delta\varkappa$, ac reliquae litterae rite
\origpage{56}
determinentur. Integrale etiam completum nulla difficultate reperietur, nam etiamsi fortasse
integrale inuentum nouam non inuoluat constantem, ponatur
\[ \frac{q\,dx}{\sqrt{Ax^{4} + 2Bx^{3} + Cxx + 2Dx + E}} = \frac{q\,du}{\sqrt{Au^{4} + 2Bu^{3} + Cuu + 2Du + E}} \]
et huius aequationis integrale completum ex antecedentibus assignare licebit; atque hinc integrale
quoque completum aequationis ex formulis disparibus constantis colligetur.

\textbf{§.~35.} Quemadmodum huius aequationis differentialis, vt a simplicissimis incipiam:
\[ \frac{dx}{\sqrt{f + gx}} = \frac{dy}{\sqrt{f + gy}} \]
integrale completum est:
\[ gg(xx + yy) - 2ggxy - 2ccg(x+y) + c^{4} - 4ccf = 0 \]
Deinde vero huius aequationis differentialis
\[ \frac{dx}{\sqrt{f + gxx}} = \frac{dy}{\sqrt{f + gyy}} \]
integrale completum est:
\[ xx + yy - 2xy\sqrt{1 + fgcc} - ccff = 0 \]
Tertio vero huius aequationis differentialis
\[ \frac{dx}{\sqrt{f + gx^{3}}} = \frac{dy}{\sqrt{f + gy^{3}}} \]
integrale completum est
\[ f(xx + yy) + \frac{ggcc}{4f}xxyy - gcxy(x+y) - 2fxy - gcc(x+y) - 2fc = 0 \]
Quarto porro huius aequationis differentialis
\[ \frac{dx}{\sqrt{f + gx^{4}}} = \frac{dy}{\sqrt{f + gy^{4}}} \]
integrale completum repertum est
\[ f(xx + yy) - fcc - gccxxyy - 2xy\sqrt{f(f + gc^{4})} = 0 \]

\origpage{57}
Ita etiam integrale completum huius aequationis
\[ \frac{dx}{\sqrt{f + gx^{6}}} = \frac{dy}{\sqrt{f + gy^{6}}} \]
reperiri poterit:

\textbf{§.~36.} Determinentur primo in §.~33. valores, ita vt prodeat haec aequatio
\[ \frac{dx}{\sqrt{fx + gx^{4}}} = \frac{dy}{\sqrt{fy + gy^{4}}} \]
cuius integralis completa reperitur:
\[ gg(xx + yy) - 4ggcxxyy - 4fgccxy(x+y) - 2ggxy - 2fgc(x+y) + ffcc = 0 \]
Ponatur nunc $x = tt$ et $y = uu$, vt prodeat haec aequatio differentialis
\[ \frac{dt}{\sqrt{f + gt^{6}}} = \frac{du}{\sqrt{f + gu^{6}}} \]
cuius propterea integralis completa erit
\[ gg(t^{4} + u^{4}) - 4ggct^{4}u^{4} - 4fgccttuu(tt+uu) - 2ggttuu - 2fgc(tt+uu) + ffcc = 0 \]
vnde notari meretur casus ex hypothesi $c = \infty$ resultans, qui dat $4gttuu(tt+uu) = f$.

\end{document}
